\begin{abstract}
Low-rank Adam methods introduce two different distortions: they omit gradient
components outside the maintained subspace, and they typically omit full-space
adaptive optimizer states for those discarded directions. \Fira-style recovery
addresses only the first distortion through a norm-based full-rank correction:
it projects the current gradient, applies a low-rank adaptive optimizer to the
projected part, and rescales the discarded component by a scalar or column-wise
ratio. We show that this mechanism is not Error Feedback. In a verified
two-dimensional stale-subspace regime, positive-$\epsilon$ recovery can drive
the projected coordinate to zero while the cumulative recovery mass remains
finite, leaving a nonzero orthogonal gradient. Lower-bounded clipping repairs
that failure, and a separated raw-gradient residual transmission mechanism,
\BECR, repairs the same regime while keeping residuals in raw-gradient units.
We then state a \BECR-SGD abstraction and evaluate the mechanism on corrected
synthetic diagnostics with explicit provenance, common random numbers, and
figure-level traceability. The resulting paper is a theory/mechanism study: it
does not claim Adam convergence, optimizer superiority, or neural validation.
\end{abstract}
